% =========================================================================
\section{Resumen operativo}\label{sec:resumen}
% =========================================================================

\begin{metodo}[{Resolución de $\Lop[y]=b$ con coeficientes constantes}]
\begin{enumerate}[label=\textcolor{ambar}{\bfseries\arabic*.}, leftmargin=1.6em]
  \item Escribir la ecuación característica $p(\lambda)=0$ y hallar sus
        raíces con su multiplicidad.
  \item Cada raíz $\lambda_k$ de multiplicidad $m_k$ aporta
        $\eul^{\lambda_kx},x\eul^{\lambda_kx},\dots,x^{m_k-1}\eul^{\lambda_kx}$; los
        pares complejos se sustituyen por sus versiones reales
        (\cref{tab:base}).
  \item Para la particular: variación de las constantes (resolver
        $W\kv'=\bv$) o función de Green (resolver $\Lop_xG=\delta$ con las
        condiciones adecuadas e integrar contra $b$); ambas son la misma
        fórmula \eqref{eq:variacion-green}.
  \item Solución general $y=y_h+y_p$; ajustar las $n$ constantes con los
        datos.
\end{enumerate}
\end{metodo}

\begin{table}[!htb]
\centering
\caption{Contribución de cada raíz del polinomio característico a la base
real de $\Ker\Lop$.}
\label{tab:base}
\renewcommand{\arraystretch}{1.5}
\small
\begin{tabular}{@{}l l c@{}}
\toprule
\textcolor{cian}{\bfseries Raíz} & \textcolor{cian}{\bfseries Soluciones} &
\textcolor{cian}{\bfseries Número}\\
\midrule
$\lambda\in\R$ simple & $\eul^{\lambda x}$ & $1$\\
$\lambda\in\R$ de multiplicidad $m$ & $x^{j}\eul^{\lambda x}$, \ $0\le j<m$ & $m$\\
$\alpha\pm\iu\beta$ simples & $\eul^{\alpha x}\cos\beta x$, \ $\eul^{\alpha x}\sin\beta x$ & $2$\\
$\alpha\pm\iu\beta$ de multiplicidad $m$ & $x^{j}\eul^{\alpha x}\cos\beta x$, \ $x^{j}\eul^{\alpha x}\sin\beta x$, \ $0\le j<m$ & $2m$\\
\bottomrule
\end{tabular}
\end{table}

\begin{observacion}[Errores típicos]\label{obs:errores}
\leavevmode
\begin{enumerate}[label=(\roman*)]
  \item Olvidar el factor $x^{j}$ en las raíces múltiples, lo que deja menos
        de $n$ soluciones.
  \item No dividir por $a_n$ en la última componente de $\bv$.
  \item Dar $\eul^{(\alpha\pm\iu\beta)x}$ como solución final cuando se pide
        una solución real.
  \item Pensar que la variación de las constantes sólo vale con
        coeficientes constantes: vale siempre que se conozca un sistema
        fundamental de la homogénea.
\end{enumerate}
\end{observacion}
